\documentclass{article}
\usepackage{amsmath}
\title{Prime Ulam Cylinder: Mathematical Foundation}
\author{Mathematica}
\begin{document}
\maketitle

\section{Introduction}

The Ulam spiral, invented by Stanislaw Ulam in 1963, reveals an
unexpected visual structure in the distribution of prime numbers. When
integers are arranged in a spiral and primes are marked, diagonal lines
emerge --- a pattern unexplained by classical number theory.

\section{The Ulam Cylinder}

The Ulam cylinder extends this spiral into three dimensions. Each
integer $n$ is mapped to coordinates $(x, y, z)$ along a helical path
of circumference $C$:
\[
  z = \left\lfloor \frac{n}{C} \right\rfloor, \quad
  \theta = \frac{2\pi n}{C}, \quad
  x = r\cos\theta, \quad y = r\sin\theta
\]
Prime integers produce a sparse but structured point cloud. Vertical
alignments in the cylinder correspond to the diagonal alignments in the
2D Ulam spiral.

\section{Prime Number Theorem}

The density of primes near $n$ approximates $\frac{1}{\ln n}$ by the
Prime Number Theorem:
\[
  \pi(n) \sim \frac{n}{\ln n}
\]
where $\pi(n)$ counts the number of primes up to $n$. The 3D
visualization makes this thinning density perceptible as vertical
sparsity in the upper cylinder.

\section{Visualization Method}

The animation renders 100 primes mapped to the cylinder surface.
Each prime $p_i$ is assigned a color interpolated by index, and the
camera orbits the structure in a full $360^\circ$ turntable over 300
frames at 4K resolution.

\end{document}
